\subsection{Linearity}

\label{sec:linearity}

An argument with multiplicity \texttt{0} is guaranteed to be erased at run
time. Correspondingly, an argument with multiplicity \texttt{1} is guaranteed
to be used exactly once. The intuition, similar to that of Linear
Haskell~\cite{linhaskell}, is that, given a function type of the form\ldots

\begin{lstlisting}
f : (1 x : a) -> b
\end{lstlisting}

\ldots then, if an expression \texttt{f x} is evaluated exactly once, 
\texttt{x} is evaluated exactly once in the process. 
%
QTT is a new core language, and the combination of linearity with dependent
types has not yet been extensively explored. Thus, we consider
the multiplicity \texttt{1} to be experimental, and in general Idris 2
programmers can get by without it---nothing in the Prelude exposes an interface
which requires it. Nevertheless, we have found
that an important application of linearity is in controlling resource
usage.
In the rest of this section, we describe two examples of this. First, 
we show in general how linearity can prevent us duplicating an argument, which can be
important if the argument represents an external resource; then, we give a 
more concrete example showing how the \TC{IO} type described in
Section~\ref{sec:interactive} is implemented.

\subsubsection{Example 1: Preventing Duplication}

To illustrate multiplicity \texttt{1} in practice, we can try (and fail!) to
write a function which duplicates a value declared as ``use once'',
interactively:

\begin{lstlisting}
dup : (1 x : a) -> (a, a)
dup x = ?dup_rhs
\end{lstlisting}

Inspecting the \texttt{dup\_rhs} hole shows that we have:

\begin{lstlisting}
 0 a : Type
 1 x : a
-------------------------------------
dup_rhs : (a, a)
\end{lstlisting}


So, \texttt{a} is not available at run-time, and \texttt{x} must be used
exactly once in the definition of \texttt{dup\_rhs}. We can write a
partial definition:

\begin{lstlisting}
dup x = (x, ?second_x)
\end{lstlisting}

However, if we check the hole \texttt{second\_x} we see that \texttt{x}
is not available, because there was only 1 and it has already been consumed:

\begin{lstlisting}
 0 a : Type
 0 x : a
-------------------------------------
second_x : a
\end{lstlisting}

We see the same result if we try \texttt{dup x = (?second\_x, x)}. If
we persist, and try\ldots

\begin{lstlisting}
dup x = (x, x)
\end{lstlisting}

\ldots then Idris reports ``\texttt{There are 2 uses of linear name x}''.

\textbf{Remark:} As we noted earlier, only usages in the body of a definition
count. This means we can still parameterise data by linear variables. For
example, if we have an \texttt{Ordered} predicate on lists, we can write
an \FN{insert} function on ordered linear lists:

\begin{lstlisting}
insert : a -> (1 xs: List a) -> (0 _ : Ordered xs) -> List a
\end{lstlisting}

The use in \TC{Ordered} does not count, and since \TC{Ordered} has multiplicity
$0$ it is erased at run time, so any occurrence of \texttt{xs} when
building the \TC{Ordered} proof also does not count.

\subsubsection{Example 2: I/O in Idris 2}

\label{sec:io}

Like Idris 1 and Haskell, Idris 2 uses a parameterised type \texttt{IO} to
describe interactive actions. Unlike the previous implementation, this is
implemented via a function which takes an abstract representation of the
outside world, of primitive type \texttt{\%World}:

\begin{lstlisting}
PrimIO : Type -> Type
PrimIO a = (1 x : %World) -> IORes a
\end{lstlisting}

The \texttt{\%World} is consumed exactly once, so
it is not possible to use previous worlds (after all, you
can't unplay a sound, or unsend a network message). It returns an
\texttt{IORes}:

\begin{lstlisting}
data IORes : Type -> Type where
     MkIORes : (result : a) -> (1 w : %World) -> IORes a
\end{lstlisting}

This is a pair of a result (with usage $\omega$), 
and an updated world state. The
intuition for multiplicities in data constructors is the same as 
in functions: here, if \texttt{MkIORes result w} is evaluated exactly once, then
the world \texttt{w} is evaluated exactly once.
We can now define \texttt{IO}:

\begin{lstlisting}
data IO : Type -> Type where
     MkIO : (1 fn : PrimIO a) -> IO a
\end{lstlisting}

There is a primitive \texttt{io\_bind} operator (which we can use to implement
$\mathtt{>>=}$), which guarantees that an action and its continuation are
executed exactly once:

\label{sec:iobind}
\begin{lstlisting}
io_bind : (1 act : IO a) -> (1 k : a -> IO b) -> IO b
io_bind (MkIO fn) = \k => MkIO (\w => let MkIORes x' w' = fn w
                                          MkIO res = k x' in res w')
\end{lstlisting}

The multiplicities of the \texttt{let} bindings are inferred from the values
being bound. Since \texttt{fn w} uses \texttt{w}, which is required to be
linear from the type of \texttt{MkIO}, \texttt{MkIORes x' w'} must itself be
linear, meaning that \texttt{w'} must also be linear. It can be informative
to insert a hole:
%to see how the multiplicities are updated:

\begin{lstlisting}
io_bind (MkIO fn)
    = \k => MkIO (\w => let MkIORes x' w' = fn w in ?io_bind_rhs)
\end{lstlisting}

This shows that, at the point \texttt{io\_bind\_rhs} is evaluated, we have
consumed \texttt{fn} and \texttt{w}, and we still have to run the continuation
\texttt{k} once, and use the updated world \texttt{w'} once:

\begin{lstlisting}
 0 b : Type
 0 a : Type
 0 fn : (1 x : %World) -> IORes a
 1 k : a -> IO b
 0 w : %World
 1 w' : %World
   x' : a
-------------------------------------
io_bind_rhs : IORes b
\end{lstlisting}

This implementation of \texttt{IO} is similar to the Haskell
approach~\cite{awkward}, with two differences:

\begin{enumerate}
\item The \texttt{\%World} token is guaranteed to be consumed exactly once, so
there is a type level guarantee that the outside world is never duplicated or
thrown away.
\item There is no built-in mechanism for exception handling, because the
type of \texttt{io\_bind} requires that the continuation is executed
exactly once. So, in \texttt{IO} primitives, we must be explicit about
where errors can occur. One can, however, implement higher level abstractions
which allow exceptions if required.
\end{enumerate}

Linearity is, therefore, fundamental to the implementation of \TC{IO} in
Idris 2. Fortunately, none of the implementation details need to be exposed
to application programmers who are using \TC{IO}.
However, once a programmer has an understanding of the combination of linear
and dependent types, they can use it to encode and verify more sophisticated
APIs. %, as we will see in the next section.
